%%% -*- coding: utf-8 -*-

\lecture{Unity Basics}{unity-basics}
\section{Unity Basics}

\showtitlepage{Scenes, Components, Lifecycle}

%%% Every example is taken from SEE, so the sources named in the listings can
%%% be opened next to the slide.

\begin{frame}{A Scene Is Assembled, Not Constructed}
  In an ordinary C\# program you decide in code which objects exist. In Unity
  most of them exist because somebody placed them in a scene and saved it:
  \texttt{Assets/Scenes/SEEStart.unity} is such a file.
  \vspace{1ex}

  \begin{itemize}
    \item A \emph{GameObject} is an empty shell: a name, a \texttt{Transform}
      and nothing else. By itself it does nothing at all.
    \item A \emph{component} is what gives it data or behaviour. A
      \texttt{Renderer} draws it, a \texttt{Collider} makes it solid, and a
      \texttt{MonoBehaviour} of your own runs code on it.
    \item A \emph{prefab} is a GameObject with its components and children
      saved as an asset, so that one scene can hold many instances of it and
      one edit can change them all.
  \end{itemize}
  \vspace{2ex}

  \hint{Nodes and edges of a dependency graph are represented as GameObjects in a code city.
      They are procedurally generated along with their necessary components.
      Their \texttt{NodeRef} or \texttt{EdgeRef} component, respectively,
      is used to refer to the underlying graph node/edge.}
\end{frame}

\begin{frame}[fragile]{Getting at the Components}
  \begin{lstlisting}[basicstyle=\scriptsize\ttfamily]
// Assets/SEE/Game/SceneManipulation/GameNodeMover.cs
Node node = child.GetComponent<NodeRef>().Value;   // null if not attached

if (gameObject.TryGetComponent(out Renderer renderer))   // the safe form
{
    renderer.enabled = false;
}

// Assets/SEE/Game/Drawable/BlinkEffect.cs
Renderer[] all = obj.GetComponentsInChildren<Renderer>();  // whole subtree
BlinkEffect effect = obj.AddComponent<BlinkEffect>();      // at run time
  \end{lstlisting}
  \vspace{1ex}

  \begin{itemize}
    \item \texttt{GetComponent} looks at one GameObject only and yields
      \texttt{null} if nothing is there; the \texttt{InChildren} and
      \texttt{InParent} variants walk the hierarchy instead.
    \item Unity overloads \texttt{==} for its \texttt{Object}: a destroyed
      object compares equal to \texttt{null} while the C\# reference still
      exists. \texttt{?.} and \texttt{??} do not know that and will happily
      call into a dead object.
  \end{itemize}
\end{frame}

\begin{exerciseframe}{A Node Prefab}
  Open \texttt{Assets/Scenes/Empty.unity} and build a miniature of what SEE
  draws:
  \vspace{1ex}

  \begin{enumerate}
    \item A cube named \texttt{Node}, carrying a component
      \texttt{NodeMarker} of your own with a label that you can type in the
      inspector without making the field public.
    \item Save the cube as a prefab and place three instances in the scene,
      each with a different label, one of them tinted red.
    \item \texttt{NodeMarker} logs its label once at start-up, and clicking
      the message in the console selects the object it came from.
  \end{enumerate}
  \vspace{1ex}

  Then change the scale of the cube in the prefab asset. Which of your three
  instances follow, and which do not? Why?
\end{exerciseframe}

\begin{solutionframe}{A Node Prefab}
  \begin{lstlisting}[basicstyle=\scriptsize\ttfamily]
public class NodeMarker : MonoBehaviour
{
    [SerializeField] private string label = "unnamed";   // visible, private

    private void Start()
    {
        bool drawn = TryGetComponent(out Renderer _);
        Debug.Log($"{name}: label={label}, drawn={drawn}", this);
    }
}
  \end{lstlisting}
  \vspace{1ex}

  \begin{itemize}
    \item \texttt{[SerializeField]} is what makes a private field appear in
      the inspector and be saved with the scene. Making it public would do
      the same, and would also let every other class write it.
    \item The second argument of \texttt{Debug.Log} is the \emph{context}
      object: the console highlights it when you click the message.
    \item All three instances grow, except where you overrode the value.
      Tinting one red overrides its material only; scale was never overridden,
      so it still comes from the asset.
  \end{itemize}
\end{solutionframe}

\begin{frame}[fragile]{The \texttt{MonoBehaviour} Lifecycle}
  \framelabel{slide:lifecycle}
  \begin{lstlisting}[basicstyle=\scriptsize\ttfamily]
Awake()        // once, as the object comes into existence
OnEnable()     // every time it becomes enabled
Start()        // once, before the first Update, only if enabled
FixedUpdate()  // per physics step: zero, one or many times per frame
Update()       // once per frame -- Time.deltaTime is how long the last took
LateUpdate()   // after every Update; cameras follow here
OnDisable()    // becomes disabled, or is about to be destroyed
OnDestroy()    // once, at the very end
  \end{lstlisting}
  \vspace{1ex}

  \begin{itemize}
    \item Every \texttt{Awake} in the scene runs before any \texttt{Start}.
      That is the rule the split exists for: find yourself in \texttt{Awake},
      reach for others in \texttt{Start}.
    \item \texttt{Awake} runs even when the component is disabled; only an
      inactive GameObject postpones it. \texttt{Start} waits for enabled.
    \item Never write a constructor. Unity creates components for you, off the
      main thread's usual rules, and a constructor cannot touch the scene.
  \end{itemize}
\end{frame}

\begin{frame}[fragile]{Coroutines}
  \begin{lstlisting}[basicstyle=\scriptsize\ttfamily]
// Assets/SEE/Game/Drawable/BlinkEffect.cs
public IEnumerator Blink()
{
    while (loopOn)                        // an ordinary loop, except that
    {                                     // it survives across frames
        renderer.enabled = false;
        yield return new WaitForSeconds(0.2f);
        renderer.enabled = true;
        yield return new WaitForSeconds(0.5f);
    }
}
...
StartCoroutine(Blink());                  // BlinkEffect does this in Start()
  \end{lstlisting}
  \vspace{0.5ex}

  \begin{itemize}
    \item A coroutine is a method returning \texttt{IEnumerator}. Each
      \texttt{yield return} hands control back: \texttt{null} resumes next
      frame, \texttt{WaitForSeconds} after scaled time,
      \texttt{WaitForSecondsRealtime} regardless of \texttt{Time.timeScale}.
    \item It is not a thread. It runs on the main thread between frames, so it
      may touch the scene freely --- and it dies with its GameObject:
      deactivating the object stops it for good, disabling only the component
      does not stop it at all.
  \end{itemize}
\end{frame}

\begin{exerciseframe}{Predict the Order}
  Two GameObjects, \texttt{A} and \texttt{B}, each carry a component that logs
  every one of \texttt{Awake}, \texttt{OnEnable}, \texttt{Start} and the first
  \texttt{Update}. The component on \texttt{A} holds a reference to \texttt{B}.
  \vspace{1ex}

  \begin{enumerate}
    \item Both are active when the scene loads. Which of the eight messages
      have a guaranteed order, and which do not?
    \item \texttt{B} is inactive when the scene loads, and \texttt{A} calls
      \texttt{b.SetActive(true)} in its \texttt{Start}. Now what is the order?
    \item \texttt{A} wants to read a value that \texttt{B} computes for
      itself. In which method does each of the two belong?
  \end{enumerate}
  \vspace{1ex}

  \hint{The lifecycle is on \slideref{slide:lifecycle}.}
\end{exerciseframe}

\begin{solutionframe}{Predict the Order}
  \begin{enumerate}
    \item \texttt{A.Awake} and \texttt{B.Awake} both precede
      \texttt{A.Start} and \texttt{B.Start}, which both precede either
      \texttt{Update}. Between \texttt{A} and \texttt{B} nothing is
      guaranteed: script execution order is arbitrary unless you configure it,
      and code that depends on it breaks when a scene is rebuilt.
    \item \texttt{A.Awake}, \texttt{A.OnEnable}, \texttt{A.Start}, then
      \texttt{B.Awake} and \texttt{B.OnEnable} inside the
      \texttt{SetActive} call, then \texttt{B.Start} --- still in the same
      frame, before any \texttt{Update}.
    \item \texttt{B} computes it in \texttt{Awake}, \texttt{A} reads it in
      \texttt{Start}. That pair is ordered by the engine, so it holds however
      the two objects happen to be visited.
  \end{enumerate}
\end{solutionframe}

\begin{exerciseframe}{Blink Without an \texttt{Update}}
  This works, but it counts time by hand and runs code every single frame to
  do something twice a second:
  \vspace{1ex}

  \begin{lstlisting}[basicstyle=\scriptsize\ttfamily]
private float elapsed;
private bool visible = true;

private void Update()
{
    elapsed += Time.deltaTime;
    if (elapsed >= 0.5f)
    {
        elapsed = 0.0f;
        visible = !visible;
        GetComponent<Renderer>().enabled = visible;
    }
}
  \end{lstlisting}
  \vspace{1ex}

  Rewrite it as a coroutine that is off for 0.2 and on for 0.5 seconds, as
  \texttt{BlinkEffect} does. Leave the object visible when the component is
  switched off.
\end{exerciseframe}

\begin{solutionframe}{Blink Without an \texttt{Update}}
  \begin{lstlisting}[basicstyle=\scriptsize\ttfamily]
private Renderer objectRenderer;
private void Start()
{
    objectRenderer = GetComponent<Renderer>();
    StartCoroutine(Blink());            // once; no per-frame work is left
}
private IEnumerator Blink()
{
    while (true)                        // the component's life bounds it
    {
        objectRenderer.enabled = false;
        yield return new WaitForSeconds(0.2f);
        objectRenderer.enabled = true;
        yield return new WaitForSeconds(0.5f);
    }
}
private void OnDisable()                // a disabled component keeps its
{                                       // coroutines running, so stop them
    StopAllCoroutines();
    objectRenderer.enabled = true;      // never leave it invisible
}
  \end{lstlisting}
\end{solutionframe}

\begin{frame}[fragile]{The Transform Hierarchy}
  \begin{lstlisting}[basicstyle=\scriptsize\ttfamily]
// Assets/SEE/Game/SceneManipulation/GameNodeMover.cs
public static void SetParent(GameObject child, GameObject newParent)
{
    child.GetComponent<NodeRef>().Value.Reparent(          // in the graph
        newParent.GetComponent<NodeRef>().Value);
    child.transform.SetParent(newParent.transform);        // in the scene
}

transform.SetParent(parent);        // keeps the world pose, rewrites local
transform.SetParent(parent, false); // keeps the local values, moves it
  \end{lstlisting}
  \vspace{1ex}

  \begin{itemize}
    \item There is no separate scene tree: the hierarchy \emph{is} the
      \texttt{Transform} hierarchy. \texttt{localPosition} is measured in the
      parent's frame, \texttt{position} in the world's, and the parent's scale
      and rotation are folded into every child below it.
    \item A parent owns its subtree. Destroying it destroys the children,
      deactivating it deactivates them --- their \texttt{activeSelf} stays
      \texttt{true} while \texttt{activeInHierarchy} turns \texttt{false}.
  \end{itemize}
\end{frame}

\begin{frame}{Prefabs, Instances, Variants}
  \begin{itemize}
    \item The \emph{asset} is the original. An \emph{instance} in a scene
      keeps a link to it, so a change to the asset reaches every instance.
    \item An instance may deviate: any value you edit on it becomes an
      \emph{override}, shown in bold in the inspector, and an override is
      immune to later changes of that same value in the asset. \emph{Apply}
      pushes it back into the asset, \emph{Revert} throws it away.
    \item A \emph{variant} is a prefab whose base is another prefab. It is the
      override idea made permanent and reusable: the variant records what
      differs, everything else keeps following the base --- including
      components added to the base afterwards.
  \end{itemize}
  \vspace{2ex}

  \hint{Copying a prefab and editing the copy looks easier and costs you the
  link. The base can then be fixed a hundred times without the copy ever
  noticing.}
\end{frame}

\begin{exerciseframe}{Reparent Without Moving}
  \texttt{parent} sits at world position $(2, 0, 0)$ with scale $2$;
  \texttt{child} sits at $(6, 0, 0)$ with scale $1$ and has no parent yet.
  \vspace{1ex}

  \begin{enumerate}
    \item After \texttt{child.transform.SetParent(parent.transform)}: what are
      the child's \texttt{position}, \texttt{localPosition} and
      \texttt{localScale}?
    \item And after \texttt{SetParent(parent.transform, false)} instead?
    \item \texttt{GameNodeMover.SetParent} does two things, not one. Which,
      and what breaks if you forget either?
    \item \texttt{Node.prefab} exists. The reflexion city needs the same node
      in red with one extra component. Variant or copy, and what happens when
      \texttt{Node.prefab} later gains a collider?
  \end{enumerate}
\end{exerciseframe}

\begin{solutionframe}{Reparent Without Moving}
  \begin{enumerate}
    \item \texttt{position} stays $(6, 0, 0)$ --- that is the point of the
      one-argument form. To keep it, Unity rewrites the local values:
      \texttt{localPosition} becomes $(2, 0, 0)$, because the child is $4$
      world units past a parent scaled by $2$, and \texttt{localScale}
      becomes $0.5$, so that the world scale is still $1$.
    \item The local values are kept instead, so \texttt{localPosition} stays
      $(6, 0, 0)$ and the child jumps to $2 + 2 \cdot 6 = (14, 0, 0)$ at
      scale $2$. This is the form for UI and for freshly created children.
    \item It reparents the graph node and the game object. Forget the graph
      and the city redraws in the old shape; forget the transform and the
      cube keeps hanging under its old neighbour.
    \item A variant. A copy would take the collider only if somebody
      remembers to add it there too.
  \end{enumerate}
\end{solutionframe}

%%% Local Variables:
%%% mode: latex
%%% TeX-master: "main-slides"
%%% mode: reftex
%%% mode: flyspell
%%% End:
